\section{Arquitetura e Padrões de Aplicações Web}
\label{sec:arquitetura-padroes-aplicacoes-web}

\subsection{Cliente-Servidor}
\label{subsec:cliente-servidor}
A arquitetura cliente-servidor organiza um sistema distribuído em provedores de
serviços e requisitantes \cite{Tanenbaum2007}. Nesse contexto, a comunicação
orientada pelo estilo REST e pelo protocolo HTTP define interações sem estado,
nas quais cada requisição é independente e deve carregar as informações
necessárias ao seu processamento \cite{rfc9110}.

\subsection{REST e HTTP}
\label{subsec:rest-http}
O estilo arquitetural REST (Representational State Transfer) define um conjunto de
restrições que, quando aplicadas, promovem escalabilidade, simplicidade e independência
entre cliente e servidor \cite{fielding2000}. A comunicação entre os componentes do
sistema ocorre por meio de recursos identificados por URIs, manipulados por métodos
HTTP (GET, POST, PUT, DELETE e etc.) e representados em formatos padronizados, como
JSON ou XML. A aplicação dessas restrições permite que o cliente e o servidor
evoluam de forma independente, desde que a interface entre eles permaneça consistente

\begin{table}[!htb]
    \centering
    \caption[Códigos HTTP e suas aplicações no módulo]{Códigos HTTP, significados padrão e aplicações no módulo.
    \label{tab:codigos-http}}
    \begin{tabular}{@{}p{2.4cm}p{4cm}p{7.7cm}@{}}
        \toprule
        \textbf{Código HTTP} & \textbf{Significado padrão} & \textbf{Aplicação no módulo} \\
        \midrule
        200 & Sucesso & Consultas e operações concluídas com êxito. \\
        201 & Recurso criado & Criação de uma nova solicitação. \\
        400 & Requisição inválida & Erros de validação nos dados enviados. \\
        401 & Falha de autenticação & Credenciais inválidas. \\
        403 & Proibido & Tentativa de acesso a recurso sem permissão. \\
        404 & Não encontrado & Solicitação ou recurso inexistente. \\
        409 & Conflito de estado & Tentativa de transição inválida (por exemplo, reverter solicitação não aprovada). \\
        \bottomrule
    \end{tabular}
    \fonte{Elaboração própria, com base em RFC 9110 \cite{rfc9110}.}
\end{table}

% TO-DO: Um diagrama de arquitetura cliente-servidor, ilustrando a comunicação 
% entre a interface, a API e o banco de dados, com destaque para a separação de 
% responsabilidades e a independência entre os componentes.


\subsection{MVC e Separação de Responsabilidade (SoC)}
\label{subsec:mvc-soc}
O padrão de arquitetura MVC (Model-View-Controller) organiza a aplicação em três
componentes: o \textit{Model} representa os dados e as regras de negócio, a \textit{View}
é responsável pela apresentação e interação com o usuário, e o \textit{Controller} atua
como intermediário, recebendo as requisições da interface, manipulando os dados e retornando
respostas \cite{gamma1994}. A separação de responsabilidades (SoC - Separation of Concerns)
permite que cada componente evolua independentemente, desde que a interface entre eles
permaneça consistente.

\subsection{Service Layer}
\label{subsec:service-layer}
O padrão Service Layer organiza a lógica de negócio em uma camada intermediária entre
o \textit{Controller} e o \textit{Model}, promovendo a reutilização de regras de negócio
e a centralização das operações que envolvem múltiplas entidades \cite{fowler2002}. Essa
abordagem facilita a manutenção do sistema, pois as alterações nas regras de negócio
podem ser implementadas em um único local, sem impactar diretamente os controladores
ou modelos individuais.

% TO-DO: Diagrama de camadas mostrando a separação entre Controller, Service Layer e Model,
% com setas indicando a direção das chamadas e a independência entre as camadas.

\subsection{Additive Change}
\label{subsec:additive-change}

O padrão Additive Change descreve a evolução de um sistema por meio da adição de novas
funcionalidades, em vez de modificar diretamente o código existente \cite{fowler2004}. Essa
abordagem minimiza o risco de introduzir erros em funcionalidades já implementadas e permite que
o sistema cresça de forma incremental, preservando a estabilidade e a confiabilidade das
operações existentes. No contexto do módulo de solicitações, o Additive Change foi aplicado ao
desenvolver novas rotinas para criação, análise e reversão de solicitações, sem alterar as
funcionalidades já presentes no Banco de Especialistas.

\subsection{OWASP e Segurança de Aplicações Web}
\label{subsec:owasp-seguranca-aplicacoes-web}
O Open Worldwide Application Security Project (OWASP) mantém o OWASP Top 10, uma lista
de referência global que elenca os dez riscos de segurança mais críticos para aplicações
web \cite{owasp2025}. A edição mais recente (2025) destaca que o Controle de Acesso Quebrado
(Broken Access Control) permanece no topo do ranking como o risco mais sério, seguido por
Configuração Insegura (Security Misconfiguration). Para um sistema que lida com dados de
usuários e decisões administrativas, como o Banco de Especialistas, alinhar o desenvolvimento
a esse referencial não é apenas uma boa prática, mas uma necessidade para mitigar
vulnerabilidades exploráveis.
